\begin{titlepage}
\thispagestyle{empty}
\begin{center}
{\large University of Rochester\\}
{\normalsize Warner School of Education\\[0.18in]}
{\color{PrimaryRed}\rule{\textwidth}{1.4pt}}\\[0.18in]
{\Huge\bfseries Many Ways to Be a Good Classmate}\\[0.18in]
{\LARGE A Kit for Kids Lesson on Peer Belonging, Communication, and Support}\\[0.22in]
{\Large EDE 448 --- Module 1}\\[0.35in]
\colorbox{SoftBlue}{\parbox{0.88\textwidth}{\centering\normalsize\textbf{From Peer Help to Peer Access}\\[0.07in]A placement-informed 10-minute peer-education lesson using OAR resources to help students practice five concrete actions: invite, wait, ask, respect, and include.}}
\end{center}
\vfill
\renewcommand{\arraystretch}{1.25}
\begin{center}
\begin{tabular}{>{\bfseries\RaggedLeft\arraybackslash}p{1.75in} >{\RaggedRight\arraybackslash}p{4.9in}}
Student & Piter Zacari Garcia Bautista \\
Assignment & Teach a Kit for Kids Lesson \\
Course & EDE 448 --- Communication and Positive Behavioral Supports for Students with Autism and Other Complex Needs \\
Primary Resources & OAR Kit for Kids, \textit{What's Up with Nick?}, and Friendship Tip Sheets \\
Teaching Context & Grade 2 technology classroom scenario grounded in non-identifying placement observations \\
Course Theme & Communication as access to dignity, agency, participation, and belonging \\
Submission Date & July 2026 \\
\end{tabular}
\end{center}
\vfill
\begin{center}
\begin{minipage}{0.9\textwidth}
\centering\large\itshape Inclusion is not only being present in the same room. It means students can communicate, participate, make choices, receive support without shame, and belong.
\end{minipage}
\end{center}
\end{titlepage}
